\SourceLeaf{318}%
für die Kugelfläche gegebenen analog sind, und in dieselben übergehen, wenn man die Excentricität \(=0\), oder die beiden Halbachsen der erzeugenden Ellipse unter sich gleich und \(=1\) setzt.

Die Breite des der conformen Übertragung auf die Kugelfläche zum Grunde liegenden Normalparallelkreises bezeichne ich (wie oben Art. 3) mit \(P\) für die Ellipsoidfläche, und mit \(Q\) für die Kugelfläche; zugleich nehme ich an, dieser Normalparallelkreis sei so gewählt, dass \(Q\) dem arithmetischen Mittel der Breiten der beiden betreffenden Punkte auf der Kugelfläche gleich wird: diese Breiten selbst seien \(Q+\frac12q\) und \(Q-\frac12q\). Es sollen ferner \(a,A,\alpha,e,\varphi,\theta\) dieselben Bedeutungen behalten, wie in der ersten Abhandlung, Art. 2. 3. 4 ff.: es bedeuten nemlich
\begin{list}{}{\setlength{\leftmargin}{1em}\setlength{\itemindent}{0pt}\setlength{\itemsep}{0pt}\setlength{\parsep}{0pt}}
\item \(a\) die halbe grosse Achse des Ellipsoids, oder den Halbmesser des Äquators,
\item \(A\) den Halbmesser der Kugel,
\item \(1:\alpha\) das constante Verhältniss der Längenunterschiede auf dem Ellipsoid zu den entsprechenden auf der Kugel,
\item \(e=\sin\varphi\) die Excentricität der erzeugenden Ellipse,
\item \(\sin\theta=e\sin P\).
\end{list}

\indent Den zwischen den beiden Punkten auf der Kugelfläche enthaltenen Grösstekreisbogen bezeichne ich mit \(s\); die Azimuthe dieses Bogens am ersten und zweiten Endpunkte mit \(U+\frac12u\) und \(U-\frac12u\pm180^0\). Erwägt man nun noch, dass der Längenunterschied zwischen beiden Punkten \(=\alpha l\) ist, so findet man zunächst die vier strengen Formeln
\[
\begin{aligned}
\sin\tfrac12s\mathord{.}\cos U&=\cos\tfrac12\alpha l\mathord{.}\sin\tfrac12q\\
\sin\tfrac12s\mathord{.}\sin U&=\sin\tfrac12\alpha l\mathord{.}\cos Q\\
\cos\tfrac12s\mathord{.}\cos\tfrac12u&=\cos\tfrac12\alpha l\mathord{.}\cos\tfrac12q\\
\cos\tfrac12s\mathord{.}\sin\tfrac12u&=\sin\tfrac12\alpha l\mathord{.}\sin Q
\end{aligned}
\]
und hieraus die näherungsweise richtigen
\[
\begin{aligned}
q&=s\mathord{.}\cos U(1+\tfrac1{24}qq-\tfrac1{24}ss+\tfrac18\alpha\alpha ll)\quad .\ .\ .\ .\ .\ .\ .\quad(1)\\
\alpha l&=s\mathord{.}\frac{\sin U}{\cos Q}(1-\tfrac1{24}ss+\tfrac1{24}\alpha\alpha ll)\quad .\ .\ .\ .\ .\ .\ .\ .\ .\ .\quad(2)\\
u&=s\mathord{.}\sin U\mathord{.}\operatorname{tang}Q(1+\tfrac1{12}ss+\tfrac1{24}uu)\quad .\ .\ .\ .\ .\ .\ .\ .\quad(3)
\end{aligned}
\]
Es ist unnöthig zu erinnern, dass in diesen drei Gleichungen \(l,q,s,u\) in Theilen des Halbmessers ausgedrückt verstanden werden. Man sieht leicht, dass sie
